% Generated from chapter-12.md - do not edit.

\chapter{Love is the Best Thought of All}%
\index{love}\nobreak%

If I had to choose just one better thought --- the best thought of all --- it would be \textbf{love}.

Love is the thought that steadies us, softens us, and brings us back to what matters most. It is the highest goodness we can offer our partners, our friends, and the people we walk through life with. When love is present, even imperfect moments feel workable. When love is expressed, relationships grow stronger almost without effort.

Love doesn't have to be grand or dramatic. Most often, love shows up through \textbf{attention} ---that complete focus with eye contact---\textbf{affection, affirmation, and appreciation}--- the small, repeatable ways we say, \emph{I see you. I care. You matter to me. You are important.}

\section{Affection: The Sweetness That Keeps Life Tasting Good}%
\index{affection}\nobreak%

Affection is the sweetness that keeps life tasting good.

Many of us grew up with very different models of affection --- or none at all. Some families hugged and touched easily; others rarely did. Some people learned how to express warmth naturally, while others never quite learned what pleased their partner or how to show love comfortably.

The good news is this: \textbf{we can learn now}.

Asking your partner what feels loving to them is not awkward --- it's a compliment.\\
``I want to show you more often how much I care. Please help me learn how.''

Even if you think you already know, ask anyway.

Some people love hugs and kisses. Others feel most loved when someone does a small chore for them, listens closely, or leaves a thoughtful note. There is no one right way to give affection --- only the willingness to pay attention.

And here's something wonderful: when you hug someone, you can't tell whether you're giving affection or receiving it. Affection flows both ways.

\section{You Never Outgrow the Need for Affection}%
\index{affection}\nobreak%

Some experts believe that more than anything else, we all crave affection, attention, and affirmation. \emph{See me. Listen to me. Hold me.}

Those needs don't disappear when we grow up.

Most of us live close to the edge of being too busy, too tired, too stressed, too frustrated. In that state, affection becomes not a luxury but a \textbf{stress reducer} --- inexpensive, effective, and immediately comforting.

Try this experiment: See if you can \textbf{overdo} warmth and kindness.\\
Let me know when that happens.

Hugs, pats, snuggles, eye contact, listening, encouragement, compliments, notes that say \emph{I love you} --- these are not trivial gestures. They are the daily nourishment of love.

\section{Affirmation: Catching Each Other Being Good}%
\index{affirmation}\nobreak%

Affirmation is love spoken out loud.

Compliments don't have to be spontaneous masterpieces or never-before-said insights to be genuine. Repeating loving truths is not boring --- it's reassuring. People forget. I forget.

You might think you shouldn't compliment unless it's completely original. Or that compliments are manipulative. But let's be honest --- we \emph{do} want something. We want connection, warmth, and closeness. That's not manipulation; that's relationship.

Notice what your partner does well.\\
Notice their humor. (And laugh at their jokes.)\\
Notice their kindness, competence, effort, and heart.

Affirmation also builds self-esteem. Many of us carry quiet doubts in a few familiar areas:

\begin{itemize}
\item Appearance --- \emph{Am I attractive?}
\item Intelligence --- \emph{Am I smart enough?}
\item Physical ability --- \emph{Am I capable or clumsy?}
\item Emotions --- \emph{Am I okay, or am I too much?}
\end{itemize}

Thoughtful affirmations help soothe those doubts:

\begin{itemize}
\item ``You look wonderful.''
\item ``You're so good at that.''
\item ``I love the way you think.''
\item ``Thank you for listening to me.''
\end{itemize}

Sometimes the most loving affirmation is simply listening to feelings, acknowledging their presence, without correcting them or trying to make them go away.

\section{Appreciation: A Habit That Pays Dividends}%
\index{appreciation}\nobreak%

Appreciation is another form of love that's hard to overdo.

The more we notice and thank our partner for what they do, the better we both feel. It's a true win-win. Gratitude encourages generosity, and kindness tends to multiply.

When I deliberately notice what he does and thank him, he feels good --- and so do I. My attention shifts away from fault-finding and toward appreciation. That shift alone is a better thought.

It does take intention. Our minds are naturally alert to problems and mistakes --- an old survival instinct, perhaps. But appreciation can be practiced.

A simple dinner-table question can change the tone of an evening:\\
``Dearest, do I thank you enough for all the wonderful things you do?''

\section{Sharing Good News --- How Do You Respond?}%
\index{good news}\nobreak%

When your partner shares good news, you have choices.

You can respond with enthusiasm and delight.\\
You can downplay it ``for their own good.''\\
You can be polite but flat.\\
Or you can ignore it altogether.

Only one of those responses builds love.

Sharing in your partner's joy --- celebrating their successes, even small ones --- strengthens goodwill and happiness. It tells them, \emph{Your happiness matters to me. Respond with enthusiasm!}

Marriage and close partnership also come with special privileges: \textbf{bragging rights}. You get to share your victories with someone who wants to be proud of you. If modesty made that difficult growing up, love gives you permission now.

\section{Thinking Loving Thoughts Changes the Relationship}%
\index{loving thoughts}\nobreak%

Here's a powerful truth: when you plan to say something loving, you must first \textbf{think loving thoughts}.

That means looking for the good, remembering what you treasure, and gently redirecting your mind away from complaints. This doesn't mean ignoring real problems --- it means choosing the right time and not letting negativity run the show.\index{negativity bias}

When I'm irritable, I sometimes pull out a little cheat sheet --- a list of why I love him. It helps me remember who he really is when my mood forgets.

You can't always force joy, but you can practice attention. Often, the emotional shift comes later --- quietly, gently --- after you persist in thinking better thoughts.

\section{The Power of Small, Frequent Acts}

Love thrives on small, frequent gestures:

\begin{itemize}
\item A long hug
\item A 20-second kiss
\item A note on the mirror
\item A quick ``I love you'' message
\item A chocolate in a desk drawer
\end{itemize}

Little things matter because they say, \emph{I'm thinking about you.}

Healthy relationships depend far more on positive interactions than on avoiding conflict. When affection, affirmation, and appreciation are plentiful, the occasional complaint lands more softly. Love creates a buffer.

\section{Can You Say ``I Love You'' Too Often?}

I don't think so.

The words stay the same, but the meaning changes with each moment of life. Sometimes \emph{I love you} is passionate. Sometimes it's reassuring. Sometimes it's a refuge from the world.

Love is not diminished by repetition. It is strengthened by it.

\section{Love as a Better Thought}%
\index{love}\nobreak%

Directing your thoughts toward love --- especially when it would be easier to complain or withdraw --- is one of the most powerful practices I know.

Love is the best thought of all.

It's not complicated. It doesn't require perfection. It simply asks for attention, warmth, and a willingness to say, again and again, in many small ways:

\emph{You matter to me.}

And what could be a better thought than that?

\textbf{What way can I send out some love today?}
